\chapter{Casos de uso}
\label{anexo:casos-uso}

Este anexo recoge los casos de uso que guían el análisis de cada iteración del Capítulo~\ref{chap:desarrollo}.
Cada caso de uso describe una interacción concreta entre un actor y el sistema, e indica la iteración en la que se aborda.
De este modo, los requisitos funcionales del sistema quedan trazados de forma explícita con su implementación.

\section{Actores}

Se identifican tres actores que interactúan con el sistema de navegación volumétrica:

\begin{itemize}
  \item \textbf{Diseñador de niveles:} configura y precalcula los volúmenes de navegación desde el editor de Unity, sin necesidad de programar. Trabaja en tiempo de diseño.
  \item \textbf{Desarrollador integrador:} escribe el código de juego que consume la \acrshort{api} del sistema para solicitar rutas y realizar consultas espaciales. Trabaja en tiempo de ejecución.
  \item \textbf{Agente:} actor secundario que representa a la entidad voladora que, en tiempo de ejecución, sigue las rutas calculadas y evita tanto la geometría como a otros agentes.
\end{itemize}

\section{Diagrama de casos de uso}

La Figura~\ref{fig:casos-uso} resume los casos de uso del sistema y su relación con cada actor.

\begin{figure}[htp]
  \centering
  \resizebox{\textwidth}{!}{%
  \begin{tikzpicture}[
      uc/.style={draw=ficblue!75!black, fill=ficblue!7, ellipse, align=center,
                 font=\scriptsize, minimum width=3.1cm, minimum height=1.05cm, inner sep=1pt},
      assoc/.style={draw=black!55},
      actorlbl/.style={font=\footnotesize\bfseries, align=center},
      persona/.pic={
        \draw[thick] (0,0) circle (3.2mm);
        \draw[thick] (0,-0.32) -- (0,-1.15);
        \draw[thick] (-0.42,-0.62) -- (0.42,-0.62);
        \draw[thick] (0,-1.15) -- (-0.34,-1.7);
        \draw[thick] (0,-1.15) -- (0.34,-1.7);
      }
    ]

    % --- Casos de uso: columna izquierda (diseño) ---
    \node[uc] (cu01) at (2.4,-1.0) {\textbf{CU-01}\\Construir el volumen};
    \node[uc] (cu02) at (2.4,-2.5) {\textbf{CU-02}\\Precalcular y persistir};
    \node[uc] (cu03) at (2.4,-4.0) {\textbf{CU-03}\\Validar integridad};
    \node[uc] (cu07) at (2.4,-5.5) {\textbf{CU-07}\\Configurar el volumen};
    \node[uc] (cu08) at (2.4,-7.0) {\textbf{CU-08}\\Depurar y ver estadísticas};

    % --- Casos de uso: columna derecha (ejecución) ---
    \node[uc] (cu04) at (6.4,-1.0) {\textbf{CU-04}\\Solicitar una ruta};
    \node[uc] (cu05) at (6.4,-2.5) {\textbf{CU-05}\\Obtener ruta suavizada};
    \node[uc] (cu06) at (6.4,-4.0) {\textbf{CU-06}\\Ruta con holgura};
    \node[uc] (cu09) at (6.4,-5.5) {\textbf{CU-09}\\Consultas espaciales};
    \node[uc] (cu10) at (6.4,-7.0) {\textbf{CU-10}\\Ruta asíncrona};
    \node[uc] (cu11) at (6.4,-8.5) {\textbf{CU-11}\\Evitar otros agentes};

    % --- Contenedor del sistema (capa de fondo) ---
    \begin{scope}[on background layer]
      \node[draw=udcgray!80!black, thick, rounded corners=4pt, fill=udcgray!8,
            fit=(cu01)(cu04)(cu08)(cu11), inner sep=12pt,
            label={[font=\bfseries]above:Sistema de navegación volumétrica}] (sys) {};
    \end{scope}

    % --- Actores ---
    \pic at (-2.8,-3.1) {persona};
    \node[actorlbl] at (-2.8,-5.1) {Diseñador\\de niveles};
    \coordinate (disenador) at (-2.8,-3.7);

    \pic at (11.2,-1.4) {persona};
    \node[actorlbl] at (11.2,-3.4) {Desarrollador\\integrador};
    \coordinate (desarrollador) at (11.2,-2.0);

    \pic at (11.2,-6.2) {persona};
    \node[actorlbl] at (11.2,-8.2) {Agente};
    \coordinate (agente) at (11.2,-6.8);

    % --- Asociaciones del diseñador de niveles ---
    \draw[assoc] (disenador) -- (cu01);
    \draw[assoc] (disenador) -- (cu02);
    \draw[assoc] (disenador) -- (cu03);
    \draw[assoc] (disenador) -- (cu07);
    \draw[assoc] (disenador) -- (cu08);

    % --- Asociaciones del desarrollador integrador ---
    \draw[assoc] (desarrollador) -- (cu04);
    \draw[assoc] (desarrollador) -- (cu05);
    \draw[assoc] (desarrollador) -- (cu09);
    \draw[assoc] (desarrollador) -- (cu10);
    \draw[assoc] (desarrollador) -- (cu11);

    % --- Asociaciones del agente ---
    \draw[assoc] (agente) -- (cu04);
    \draw[assoc] (agente) -- (cu06);
    \draw[assoc] (agente) -- (cu11);

  \end{tikzpicture}%
  }
  \caption[Diagrama de casos de uso del sistema de navegación volumétrica.]
  {Diagrama de casos de uso del sistema de navegación volumétrica.
  El diseñador de niveles interactúa con el sistema en tiempo de diseño, mientras que el desarrollador integrador y el agente lo hacen en tiempo de ejecución.}
  \label{fig:casos-uso}
\end{figure}

\section{Descripción de los casos de uso}

A continuación se detalla cada caso de uso, indicando sus actores, su descripción, sus condiciones y la iteración del Capítulo~\ref{chap:desarrollo} en la que se aborda.

\medskip
\noindent\begin{tabular}{@{}>{\bfseries}p{2.8cm}p{9.9cm}@{}}
\rowcolor{udcpink!25}
\multicolumn{2}{@{}p{12.9cm}@{}}{\textbf{CU-01 — Construir el volumen de navegación}}\\
Actores & Diseñador de niveles. \\
Descripción & A partir de la geometría de una escena, el sistema genera el \acrshort{svo} que representa su espacio libre navegable. \\
Precondición & La escena contiene la geometría (\textit{colliders}) sobre la que se quiere navegar. \\
Postcondición & Existe en memoria un \acrshort{svo} conectado que representa el espacio navegable de la escena. \\
Iteración & 1 (tarea T3). \\
\end{tabular}

\medskip
\noindent\begin{tabular}{@{}>{\bfseries}p{2.8cm}p{9.9cm}@{}}
\rowcolor{udcpink!25}
\multicolumn{2}{@{}p{12.9cm}@{}}{\textbf{CU-02 — Precalcular y persistir el volumen (\textit{bake})}}\\
Actores & Diseñador de niveles. \\
Descripción & El diseñador ejecuta el precálculo del volumen desde el editor y lo almacena en disco para evitar reconstruirlo en cada ejecución. \\
Precondición & La geometría de la escena está definida. \\
Postcondición & El volumen queda guardado en un \textit{asset}, acompañado de un \textit{hash} de la geometría. \\
Iteración & 1 (tarea T3). \\
\end{tabular}

\medskip
\noindent\begin{tabular}{@{}>{\bfseries}p{2.8cm}p{9.9cm}@{}}
\rowcolor{udcpink!25}
\multicolumn{2}{@{}p{12.9cm}@{}}{\textbf{CU-03 — Detectar un volumen precalculado obsoleto}}\\
Actores & Sistema (al cargar la escena), en beneficio del diseñador de niveles. \\
Descripción & Al cargar el nivel, el sistema compara el \textit{hash} de la geometría actual con el almacenado y avisa si el volumen precalculado ha quedado obsoleto. \\
Precondición & Existe un volumen precalculado asociado a la escena. \\
Postcondición & El sistema confirma que el volumen sigue siendo válido o señala que debe reconstruirse. \\
Iteración & 1 (tarea T3). \\
\end{tabular}

\medskip
\noindent\begin{tabular}{@{}>{\bfseries}p{2.8cm}p{9.9cm}@{}}
\rowcolor{udcpink!25}
\multicolumn{2}{@{}p{12.9cm}@{}}{\textbf{CU-04 — Solicitar una ruta navegable entre dos puntos}}\\
Actores & Desarrollador integrador, Agente. \\
Descripción & Dado un punto de origen y uno de destino, el sistema calcula una trayectoria libre de colisiones a través del espacio navegable. \\
Precondición & Existe un volumen de navegación cargado. \\
Postcondición & Se devuelve una ruta válida o un fallo controlado si no existe ninguna. \\
Iteración & 2 (tarea T4). \\
\end{tabular}

\medskip
\noindent\begin{tabular}{@{}>{\bfseries}p{2.8cm}p{9.9cm}@{}}
\rowcolor{udcpink!25}
\multicolumn{2}{@{}p{12.9cm}@{}}{\textbf{CU-05 — Obtener una trayectoria suavizada y libre de colisiones}}\\
Actores & Desarrollador integrador, Agente. \\
Descripción & La ruta calculada se postprocesa (atajo por línea de visión y suavizado por \textit{spline}) para lograr un movimiento natural sin perder seguridad. \\
Precondición & Se ha calculado una ruta en bruto (CU-04). \\
Postcondición & Se devuelve una trayectoria suavizada que permanece en espacio libre. \\
Iteración & 2 (tarea T4). \\
\end{tabular}

\medskip
\noindent\begin{tabular}{@{}>{\bfseries}p{2.8cm}p{9.9cm}@{}}
\rowcolor{udcpink!25}
\multicolumn{2}{@{}p{12.9cm}@{}}{\textbf{CU-06 — Generar rutas con holgura para el tamaño del agente}}\\
Actores & Desarrollador integrador, Agente. \\
Descripción & El sistema considera el radio del agente al construir el volumen, de modo que las rutas conserven la holgura necesaria para no rozar la geometría. \\
Precondición & Se conoce el radio del agente en el momento de construir el volumen. \\
Postcondición & Toda ruta marcada como libre garantiza, por construcción, el margen necesario para el tamaño del agente. \\
Iteración & 3 (extensión de la tarea T4). \\
\end{tabular}

\medskip
\noindent\begin{tabular}{@{}>{\bfseries}p{2.8cm}p{9.9cm}@{}}
\rowcolor{udcpink!25}
\multicolumn{2}{@{}p{12.9cm}@{}}{\textbf{CU-07 — Configurar el volumen desde el editor de Unity}}\\
Actores & Diseñador de niveles. \\
Descripción & Desde el inspector de Unity, el usuario configura el volumen (modo de construcción y parámetros) sin escribir código. \\
Precondición & El componente \texttt{NavVolumeSpace} está añadido a la escena. \\
Postcondición & El volumen queda configurado según los parámetros elegidos. \\
Iteración & 4 (tarea T5). \\
\end{tabular}

\medskip
\noindent\begin{tabular}{@{}>{\bfseries}p{2.8cm}p{9.9cm}@{}}
\rowcolor{udcpink!25}
\multicolumn{2}{@{}p{12.9cm}@{}}{\textbf{CU-08 — Depurar el volumen y consultar sus estadísticas}}\\
Actores & Diseñador de niveles. \\
Descripción & El usuario inspecciona visualmente el volumen mediante Gizmos y consulta estadísticas como la memoria estimada o el porcentaje de vóxeles ahorrados. \\
Precondición & Existe un volumen construido. \\
Postcondición & El usuario obtiene información visual y numérica para ajustar la configuración. \\
Iteración & 4 (tarea T5). \\
\end{tabular}

\medskip
\noindent\begin{tabular}{@{}>{\bfseries}p{2.8cm}p{9.9cm}@{}}
\rowcolor{udcpink!25}
\multicolumn{2}{@{}p{12.9cm}@{}}{\textbf{CU-09 — Realizar consultas espaciales sobre el volumen}}\\
Actores & Desarrollador integrador. \\
Descripción & El sistema responde a consultas como comprobar si un punto es navegable, ajustarlo al \gls{voxel} libre más cercano o generar una posición aleatoria válida. \\
Precondición & Existe un volumen de navegación cargado. \\
Postcondición & Se devuelve el resultado de la consulta espacial solicitada. \\
Iteración & 4 (tarea T5). \\
\end{tabular}

\medskip
\noindent\begin{tabular}{@{}>{\bfseries}p{2.8cm}p{9.9cm}@{}}
\rowcolor{udcpink!25}
\multicolumn{2}{@{}p{12.9cm}@{}}{\textbf{CU-10 — Solicitar rutas de forma asíncrona}}\\
Actores & Desarrollador integrador. \\
Descripción & El sistema calcula rutas en un hilo secundario para no bloquear el juego, y permite cancelar la petición en curso. \\
Precondición & Existe un volumen de navegación cargado. \\
Postcondición & Se entrega la ruta sin bloquear el hilo principal, o bien se cancela la operación. \\
Iteración & 4 (tarea T5). \\
\end{tabular}

\medskip
\noindent\begin{tabular}{@{}>{\bfseries}p{2.8cm}p{9.9cm}@{}}
\rowcolor{udcpink!25}
\multicolumn{2}{@{}p{12.9cm}@{}}{\textbf{CU-11 — Evitar colisiones entre agentes en movimiento}}\\
Actores & Agente, Desarrollador integrador. \\
Descripción & En tiempo de ejecución, cada agente ajusta su velocidad para evitar a los demás agentes y a la geometría cercana mientras sigue su ruta. \\
Precondición & Varios agentes se desplazan simultáneamente por el escenario. \\
Postcondición & Los agentes se mueven sin colisionar entre sí ni con la geometría estática. \\
Iteración & 6 (tarea T8). \\
\end{tabular}
